\documentclass{article}
\usepackage{hyperref}
\usepackage[a4paper]{geometry}
\usepackage{fancyhdr}
\pagestyle{fancy}
\lhead{Interferenzmuster mit Quanten}
\rhead{März 2026}
\begin{document}
\section{Interferenzmuster mit Quanten} 
Wird ein Experiment, bei dem ein Interferenzmuster entsteht, durchgeführt, wie das \hyperref[Das Doppelspaltexperiment]{Doppelspaltexperiment}, nur dass die Lichtphotonen durch Quanten, wie Elektronen erstetzt werden (und das Experiment so umgebaut wird, dass es die Quanten messen kann, z.\,B. dass der Schirm durch einen ortsauflösenden Detektor ersetzt wird), so bleibt das Ergebnis das Gleiche. 
 
\subsection{Stochastik}  
Dies kann dahin weitergeführt werden, dass die Quanten nicht einmal zur gleichen Zeit das Experiment durchlaufen müssen; auch wenn zu jedem Zeitpunkt höchstens ein Elektron das Doppelspaltexperiment durchläuft bildet sich ein Interferenzmuster.
 
Der Auftreffort von einzelnen Elektronen kann dabei nicht genau vorausgesagt werden. Es kann lediglich angemerkt werden, dass die Wahrscheinlichkeit durch die Wahrscheinlichkeitsverteilung, welche das Quadrat der Amplitude der Sinuskurve des Ortes, darstellt und dass der gleichen Verteilung nach mit genug Elektronen ein Interferenzmuster entsteht.
\end{document}